\section*{第 \ref{ch:b3:cancer-biology} 章——癌症生物学}

\begin{solution}{exo:b3:cancer-biology:1}
\omterm{def:b3:cancer-biology:genes}{癌基因}是促进增殖或存活的基因的功能获得型形式，
只需一个突变等位基因即可（显性）：
\emph{RAS}（点突变）、\emph{MYC}（易位、扩增），
\emph{HER2}、\emph{BCR--ABL}。抑\omterm{def:b3:cancer-biology:genes}{癌基因}抑制
增殖，或强制细胞停滞、死亡或修复；两个等位基因都必须
丢失（细胞水平上为隐性）：\emph{RB}、\emph{TP53}、\emph{APC}、
\emph{BRCA1}。
\end{solution}

\begin{solution}{exo:b3:cancer-biology:2}
持续增殖：\emph{RAS}、\emph{MYC}。对生长抑制
因子不敏感：\emph{RB}、\emph{CDKN2A}（p16）。抵抗死亡：
\emph{BCL2}、\emph{TP53}。复制性永生：\omterm{def:b3:cancer-biology:telomerase}{端粒酶}
启动子。\omterm{prop:b3:cancer-biology:oxygen}{血管生成}：经 HIF 诱导 \emph{VEGF}（\emph{VHL}
丢失）。\omterm{def:b3:cancer-biology:metastasis}{侵袭}与\omterm{def:b3:cancer-biology:hallmarks}{转移}：E-cadherin 丢失、
\omterm{def:b3:cancer-biology:metastasis}{上皮–间质转化}相关因子。基因组不稳定：\omterm{prop:b3:dna-repair:fidelity}{错配修复}
基因、\emph{\omterm{def:b3:dna-repair:dsb}{BRCA}}。代谢失调：\emph{MYC}、PI3K 通路。
\omterm{prop:b3:cancer-biology:immune}{免疫逃逸}：\emph{PD-L1}、MHC 丢失。促\omterm{def:b3:cancer-biology:hallmarks}{肿瘤}炎症
\emph{NF-$\kappa$B} 信号。
\end{solution}

\begin{solution}{exo:b3:cancer-biology:3}
有家族史的患儿双眼都出现了多个\omterm{def:b3:cancer-biology:hallmarks}{肿瘤}，
且发病较早；散发病例则仅出现单个\omterm{def:b3:cancer-biology:hallmarks}{肿瘤}，发病较晚。他由此推断，
细胞中需要两次事件：家族性患者已遗传了其中一次，
只需再发生一次体细胞事件（其发生频率足以在多个视网膜细胞中出现），
而散发病例则需要两次体细胞事件同时发生在同一个细胞中——
罕见、发病晚、只出现单个\omterm{def:b3:cancer-biology:hallmarks}{肿瘤}。由此可知，两个拷贝都须丢失的基因
即第一个抑\omterm{def:b3:cancer-biology:genes}{癌基因}。
\end{solution}

\begin{solution}{exo:b3:cancer-biology:4}
将需要组氨酸的 \emph{Salmonella} 突变株接种于
不含组氨酸并加有受试化合物的平板上。仅那些
发生新突变、使缺陷得以回复的细胞才能形成菌落，因此菌落数即反映
其致突变性。加入肝脏提取物，是因为许多\omterm{prop:b3:cancer-biology:causes}{致癌物}
在肝酶作用下转化为活性代谢物之前并无害，
体内正是如此；而细菌缺乏这些酶。
\end{solution}

\begin{solution}{exo:b3:cancer-biology:5}
年龄每翻一倍，发病率上升 $100$ 倍：$2^{k-1} = 100$，$k - 1 =
\log_{2}100 = 6.6$，因此约需七至八个限速事件，若按
最简单的理解；若允许两次打击之间发生克隆扩增，则所需事件更少。
\end{solution}

\begin{solution}{exo:b3:cancer-biology:6}
$L = \sqrt{2Dc_{0}/q}$：当 $q = 0.01$ 时，$\sqrt{2\times 2\times 10^{-9}
\times 0.05/0.01} = \qty{141}{\micro m}$；当 $q = 0.03$ 时，
\qty{82}{\micro m}。生长迅速的\omterm{def:b3:cancer-biology:hallmarks}{肿瘤}呼吸更快，氧气在
更短距离内即被耗尽，其中心便会在较小尺寸时死亡。
\end{solution}

\begin{solution}{exo:b3:cancer-biology:7}
$10^{11}\times 10^{-3n} < 1$ 要求 $n > 3.7$：需四个疗程。两个疗程之后
余下 $10^{5}$ 个细胞——仅约十分之一立方毫米，远
低于可被检出的 $10^{9}$ 个细胞，因此经各项检测，\omterm{def:b3:cancer-biology:hallmarks}{肿瘤}
均已“消失”，而仍有十万个细胞残存；治疗
必须按计划疗程数持续进行，不应以\omterm{def:b3:cancer-biology:hallmarks}{肿瘤}
消失为终点。
\end{solution}

\begin{solution}{exo:b3:cancer-biology:8}
两种\omterm{def:b3:dna-repair:lesions}{损伤}都激活同一条通路（受体 $\to$ \omterm{def:b3:cancer-biology:genes}{Ras} $\to$ ERK）；
一旦其中一种存在，另一种便不再带来额外优势，因而不会被
选择。EGFR 抑制剂阻断受体；获得
下游 KRAS 突变的细胞，在受体仍被阻断的情况下恢复了信号，
因而在治疗期间被选择——这是一种通过旁路产生的抗性。
\end{solution}

\begin{solution}{exo:b3:cancer-biology:9}
E7 使 \omterm{prop:b3:cell-cycle-apoptosis:rb}{E2F} 从 Rb 上释放，驱动被感染的上皮细胞越过
\omterm{def:b3:cell-cycle-apoptosis:checkpoints}{限制点}，而无需生长因子；E6 使 \omterm{def:b3:cancer-biology:genes}{p53} 降解，因此
这种不当增殖及其造成的\omterm{def:b3:dna-repair:lesions}{损伤}既不会引发
细胞停滞，也不会引发\omterm{def:b3:cell-cycle-apoptosis:apoptosis}{细胞凋亡}。两项\omterm{def:b3:cancer-biology:hallmarks}{癌症}特征同时被提供，但
其余的——永生化、\omterm{def:b3:cancer-biology:metastasis}{侵袭}和不稳定——仍需进一步的
体细胞突变，这些突变需要数十年才能在持续感染的
细胞中积累。疫苗可诱导产生针对衣壳的中和抗体，
从而阻止感染；若在接触前接种，
则永远不会出现可被转化的感染细胞。
\end{solution}

\begin{solution}{exo:b3:cancer-biology:10}
第 $i$ 次打击之后，克隆扩大 $g$ 倍，因此处于下一次
打击风险中的细胞数为原来的 $g^{i}$ 倍；单个原始细胞
繁衍的谱系中出现完整序列的概率为 $g\cdot g^{2}\cdots$
——若每次打击的 $g$ 为常数，则累积风险为 $N(ut)^{k}
g^{k-1}$（相差与次序有关的因子），发病率仍保持
$t^{k-1}$ 的幂次，只是前置因子增大 $g^{k-1}$ 倍。若
克隆扩增本身随时间增长（呈指数增长的克隆），则曲线
进一步变陡，因此斜率 $6$ 可由少于七个
驱动因素产生：阿米蒂奇–多尔的 $k$ 是限速驱动因素数目的上界，
而经测序的\omterm{def:b3:cancer-biology:hallmarks}{肿瘤}显示为两到八个。
\end{solution}

\begin{solution}{exo:b3:cancer-biology:11}
(a) 单用药物 A 时，面对的是 $\mu_{A}N = 100$ 个预先存在的耐药细胞：
$P_{0} = e^{-100} \approx 0$；对 A 耐药的克隆得以存活，
并在 A 用药的十八周内重新增殖，到 B 开始时，它已
再次成为大种群，因此 B 同样面对 $\mu_{B}N' \gg 1$——
该序列两次失败。(b) 联合用药时，双重耐药细胞产生的概率为
每次分裂 $10^{-14}$，$\mu_{A}\mu_{B}N = 10^{-5}$，$P_{0} =
0.99999$。自始即联合用药（此时 $N$ 最小）是应遵循的
原则，因为两个小速率的乘积使抗性
难以出现。
\end{solution}

\begin{solution}{exo:b3:cancer-biology:12}
细胞数增至 $100$ 倍而 $u$ 与 $k$ 不变，则累积风险
$N(ut)^{k}$ 将高出 $100$ 倍——每头大象都应在年轻时死于
\omterm{def:b3:cancer-biology:hallmarks}{癌症}。事实并非如此，因此 $u$ 或 $k$ 必定不同：二十份
\emph{TP53} 拷贝使细胞对\omterm{def:b3:dna-repair:lesions}{损伤}的凋亡反应更强，在突变细胞扩增之前
将其清除（使该步的有效 $u$ 降低）；大型动物的单细胞突变率也更低，
原因在于修复更好，且细胞更新较慢、每个细胞每年分裂次数更少，
并且可能需要更多次打击，
即多一层抑癌机制。
抑癌能力随体型一同演化。
\end{solution}

\begin{solution}{pb:b3:cancer-biology:1}
\textbf{1.} $10^{12}/10^{3} = 10^{9}$ 个细胞；$\log_{2}10^{9} \approx 30$
次倍增。
\textbf{2.} $30\times 100 = 3000$ 天，约 $8$ 年。
\textbf{3.} \qty{1}{kg} $\approx 10^{12}$ 个细胞：再增加十次倍增，
$1000$ 天，不到三年——约为发现所需时间的三分之一。
\textbf{4.} 早期倍增时间必须约为 $60$ 天，否则\omterm{def:b3:cancer-biology:hallmarks}{肿瘤}
在体积小时生长更快，随体积增大而减慢（常见情形，
因供应与空间受限）。
\textbf{5.} 周期为 $2$ 天时，种群每 $2$
天翻倍；倍增时间为 $100$ 天，则净增长仅为出生率的 $2\,\%$：
几乎每个新生细胞都与一个死亡细胞相抵，或处于循环中的细胞很少——
出生与死亡几乎平衡，
增长只是其中的微小差值。
\textbf{6.} 药物杀伤处于循环中的细胞，因此循环细胞比例小的\omterm{def:b3:cancer-biology:hallmarks}{肿瘤}
每疗程损失的比例也较小；但快速增殖的\omterm{def:b3:cancer-biology:hallmarks}{肿瘤}损失更多，
也会在疗程之间重新生长。总的来看，
化疗能够治愈的正是快速增殖的\omterm{def:b3:cancer-biology:hallmarks}{肿瘤}
（白血病、淋巴瘤、睾丸癌）。
\textbf{7.} 比值 $300$ 除以年龄之比 $80/30 = 2.67$：$k - 1 =
\ln 300/\ln 2.67 = 5.8$，$k \approx 7$。
\textbf{8.} $N(ut)^{k} = 10^{8}\times (7\times 10^{-5})^{7} \approx
10^{-21}$：与约 $1$ 比 $3$ 的一生风险相比荒谬。缺失的因素包括：
每次打击后克隆扩增使处于风险中的细胞倍增，
不稳定性提高突变率，分裂次数多得多（干细胞循环数十年），
以及真正的\omterm{prop:b3:cancer-biology:multistep}{驱动突变}更少。
\textbf{9.} 乘以 $g^{k-1} = 100^{6} = 10^{12}$ 得 $10^{-9}$——
七次事件时仍然极小；若只有四次事件并计入扩增，
则 $10^{8}\times (7\times 10^{-5})^{4}\times 10^{6} \approx 2\times
10^{-3}$，量级正确。
真实图景是\omterm{prop:b3:cancer-biology:multistep}{驱动突变}少、克隆扩增大。
\textbf{10.} $(ut)^{k-1}/(ut)^{k} = 1/(ut) = 1/(10^{-6}\times 40) =
2.5\times 10^{4}$：携带者四十岁时的发病率
是散发病例的数万倍——遗传性\omterm{def:b3:cancer-biology:hallmarks}{癌症}发病年龄较轻。
\textbf{11.} 一次事件快十倍：发病率 $\times 10$；两次：
$\times 100$。观测到的十五至二十五倍风险表明，吸烟作用于
不止一个步骤。
\textbf{12.} 戒烟后，对尚未发生的事件而言 $u$ 回落，
因此发病率不再陡峭上升；但细胞中已累积的打击依然存在，
戒烟者的细胞因而比从未吸烟者的细胞更接近癌变，
超额风险持续存在，
是被冻结而非被抹除。
\textbf{13.} $L = \sqrt{2\times 2\times 10^{-9}\times 0.05/0.03} =
\qty{82}{\micro m}$。
\textbf{14.} 直径不超过 $2L \approx
\qty{160}{\micro m}$ 时充分供氧。直径为 \qty{2}{mm} 的球体，其存活外缘厚
\qty{82}{\micro m}：存活比例为 $1 - (918/1000)^{3} = 0.23$。
\textbf{15.} 居中的细胞距血管 \qty{75}{\micro m}，恰在 $L$ 之内——
原则上可获得氧气，但\omterm{def:b3:cancer-biology:hallmarks}{肿瘤}血管输送的氧气较少且血流间歇，
因此居中的细胞处于缺氧状态。缺氧细胞对放射线的抗性
高两至三倍（\omterm{def:b3:dna-repair:lesions}{损伤}要成为永久性的，需要氧气参与），
而缺氧会引发依赖 \omterm{def:b3:cancer-biology:genes}{p53} 的凋亡，
因此它会选择那些丢失了 \omterm{def:b3:cancer-biology:genes}{p53} 的细胞。
\textbf{16.} 密度减半使血管间距扩大为 $\sqrt{2}$ 倍：
血管相距 \qty{212}{\micro m}，居中的细胞距血管 \qty{106}{\micro m}，
超出 $L$：它们或死亡，或在缺氧中存活，
并变得更具\omterm{def:b3:cancer-biology:metastasis}{侵袭}性、促\omterm{prop:b3:cancer-biology:oxygen}{血管生成}。
\textbf{17.} 需要 $30/2 = 15$ 倍的葡萄糖。因此\omterm{def:b3:cancer-biology:hallmarks}{肿瘤}
会富集葡萄糖类似物，
正电子发射的氟代脱氧葡萄糖扫描可使其显影。
\textbf{18.} 缺乏营养的\omterm{def:b3:cancer-biology:hallmarks}{肿瘤}会缩小至受扩散限制的大小，
并保持休眠和缺氧状态，或借用已有血管；
缺氧会选择\omterm{def:b3:cancer-biology:metastasis}{侵袭}性克隆。但修剪混乱的\omterm{def:b3:cancer-biology:hallmarks}{肿瘤}血管可降低间质压力、
使血流趋于均匀，因此化疗药物渗透更好——
两者联合有效，
而单用任一种则无效。
\textbf{19.} $10^{9}\times 10^{-2n} < 1$：$n > 4.5$，即五个疗程。
\textbf{20.} $21$ 天 $= 0.21$ 次倍增，再生长 $\times 1.16$；
每个周期的净因子为 $0.0116$；$\ln 10^{-9}/\ln 0.0116 = 4.65$：
仍需五个周期。
\textbf{21.} 期望的耐药细胞数为 $10^{-7}\times 10^{9} = 100$；
$P_{0} = e^{-100} \approx 0$：耐药必然出现。
\textbf{22.} $\mu_{1}\mu_{2}N = 10^{-14}\times 10^{9} = 10^{-5}$；
$P_{0} = 0.99999$。当细胞数为 $10^{12}$ 时，为 $10^{-2}$，$P_{0} = 0.99$。
\textbf{23.} 细胞数为 $10^{6}$：单药 $\mu N = 0.1$，$P_{0} = 0.90$；双药
时为 $10^{-8}$，$P_{0} \approx 1$。应在种群较小时治疗，
并采用联合用药。
\textbf{24.} T~细胞同时攻击大量\omterm{prop:b3:cancer-biology:immune}{新抗原}，因此没有任何单一突变能够逃脱；
逃逸需要丢失抗原呈递本身（MHC、$\beta_{2}$-微球蛋白），
或抑制 T~细胞，这是另一类更为罕见的事件。
同样高的突变负荷使\omterm{def:b3:cancer-biology:hallmarks}{肿瘤}必然对激酶抑制剂产生抗性，
也使其拥有大量\omterm{prop:b3:cancer-biology:immune}{新抗原}：
\omterm{prop:b3:dna-repair:fidelity}{错配修复}缺陷、紫外线和吸烟诱发的\omterm{def:b3:cancer-biology:hallmarks}{肿瘤}
对治疗的反应最好。
\textbf{25.} 约 $8$ 年后发现；存活\omterm{def:b3:cancer-biology:hallmarks}{肿瘤}在血管周围深 \qty{82}{\micro m}
处；双药联合时，
在 \qty{1}{cm^3} 的\omterm{def:b3:cancer-biology:hallmarks}{肿瘤}中不遇到耐药细胞的概率为 $P_{0} = 0.99999$。
\end{solution}
